\chapter{Harmonic analysis and spectral time series}
A continuous signal is a function $x:\R\to\R$ or $\C$.
The Fourier transform, when meaningful in the appropriate sense, is
$X(f)=\int_{\R}x(t)e^{-2\pi ift}\,dt$.
A length-$N$ recording is a different object: a discrete sequence.
An FFT of the recorded sequence does not directly reveal a unique
continuous generating signal without sampling assumptions.

\entry{57}{Euler's formula and Fourier orthogonality}
{$e^{i\theta}=\cos\theta+i\sin\theta$ and
$\int_0^T e^{2\pi i(k-\ell)t/T}\,dt=0$ for distinct integers
$k,\ell$.}
{Derive discrete spectral coefficients and phase representations.}
\entry{58}{Parseval's identity and Plancherel's theorem}
{For appropriate Fourier expansions/transforms,
energy in time equals energy in frequency, e.g.\
$\|x\|_2^2=\|X\|_2^2$ under a unitary convention.}
{Normalize power and check spectral energy.}
\entry{59}{Fourier convolution theorem}
{For suitable integrable functions, $\Fourier(f*g)=
\Fourier(f)\Fourier(g)$, with conventions fixed.}
{Relate a digital filter's impulse response to its transfer function.}
\entry{60}{Fourier shift and modulation identities}
{Translation of a signal multiplies its spectrum by a complex
phase; multiplication by a complex sinusoid translates frequency.}
{Understand lags, phase offsets, and analytic demodulation.}
\entry{61}{Heisenberg--Gabor uncertainty inequality}
{With standard centering and Fourier conventions, a nonzero
finite-energy waveform has time spread and frequency spread
whose product obeys a nonzero lower bound.}
{Understand the temporal and spectral resolution trade-off.}
\entry{62}{Paley--Wiener theorems}
{Appropriate compact-support and exponential-type conditions
characterize special classes of Fourier transforms;
compactly supported time-domain signals have entire Fourier
transforms of controlled exponential growth under standard hypotheses.
In particular, nonzero ordinary signals cannot have both time
and frequency support compact: the identity theorem supplies
this corollary.}
{Explain the role of analyticity in exact localization;
Benedicks-type results strengthen restrictions for finite-measure supports.}
\entry{63}{Riemann--Lebesgue lemma}
{The Fourier transform of an $L^1$ function vanishes
at infinity.}
{Understand high-frequency decay under integrability assumptions.}
\entry{64}{Bochner and Herglotz spectral representation}
{Continuous positive-definite functions on $\R$ are Fourier
transforms of finite nonnegative measures; discrete-time analogues
use nonnegative measures on the frequency circle.}
{Characterize admissible stationary covariance functions and spectra.}
\entry{65}{Cramer's spectral representation}
{A zero-mean weakly stationary second-order process can be
represented via orthogonal-increment random spectral measures
under the standard hypotheses.}
{Interpret stochastic LFP spectra rather than deterministic FFTs alone.}
\entry{66}{Wold decomposition}
{A purely nondeterministic covariance-stationary discrete-time
process has a one-sided linear innovations representation;
general stationary processes decompose into deterministic and
purely nondeterministic parts.}
{Understand assumptions behind linear innovations and VAR models.}
\entry{67}{Yule--Walker equations}
{For a stationary autoregressive process, lagged covariance
satisfies linear equations determined by AR coefficients.}
{Relate autocovariance to AR spectral estimators.}
\entry{68}{Wiener--Hopf normal equations}
{The least-squares optimal linear predictor has a residual
orthogonal to the predictor space, yielding covariance normal equations.}
{Derive linear decoding and linear prediction filters.}
\entry{69}{Spectral coherence bound}
{A positive-semidefinite two-signal spectral matrix satisfies
$|S_{xy}(f)|^2\leq S_{xx}(f)S_{yy}(f)$ where the spectra exist;
hence $0\leq C_{xy}(f)\leq1$ for positive denominators.}
{Interpret magnitude-squared LFP coherence.}
\entry{70}{Hilbert-transform Fourier multiplier identity}
{On $L^2(\R)$, with appropriate Fourier convention,
$\Fourier(\Hilbert x)(f)=-i\operatorname{sgn}(f)\Fourier x(f)$.}
{Build analytic signals and interpret instantaneous phase.}

\begin{definition}
For zero-mean jointly wide-sense stationary signals $x,y$,
let $R_{xy}(\tau)=\E[x(t+\tau)\overline{y(t)}]$.
Under appropriate integrability, $S_{xy}=\Fourier R_{xy}$.
Magnitude-squared coherence is
$C_{xy}=|S_{xy}|^2/(S_{xx}S_{yy})$ wherever the
denominator is strictly positive.
\end{definition}
\begin{caution}
Coherence is spectral dependence, not anatomical connectivity.
Common reference, volume conduction, leakage, sample-size bias,
and nonstationarity can distort interpretation
\cite{bastos2016,cohen2018}.
\end{caution}
% --- expanded edition (October 2026) ---
\section{Proofs of selected harmonic-analysis results}
\begin{proof}[Proof of Result 59 for integrable functions]
For $f,g\in L^1(\R)$, Tonelli's theorem gives
\[
\int\!\!\int|f(s)||g(t-s)|\,ds\,dt=\|f\|_1\|g\|_1<\infty,
\]
so $f*g\in L^1$ and Fubini allows the exchange
\[
\Fourier(f*g)(\xi)=\int\!\!\int f(s)g(t-s)e^{-2\pi i\xi t}\,ds\,dt
=\int f(s)e^{-2\pi i\xi s}\Bigl(\int g(u)e^{-2\pi i\xi u}\,du\Bigr)ds,
\]
after the substitution $u=t-s$. The inner integral is $\Fourier g(\xi)$.
\end{proof}
\begin{proof}[Proof of Result 63 (Riemann--Lebesgue)]
For $f=\ind_{[a,b]}$, $|\Fourier f(\xi)|=|e^{-2\pi i\xi a}-e^{-2\pi i\xi b}|/(2\pi|\xi|)\le1/(\pi|\xi|)\to0$.
By linearity the claim holds for step functions. Every $f\in L^1$ is an
$L^1$-limit of step functions $f_n$, and
$\sup_\xi|\Fourier f(\xi)-\Fourier f_n(\xi)|\le\|f-f_n\|_1$, so the limit
property passes to $f$.
\end{proof}
The rate of decay reflects smoothness: a jump gives $1/|\xi|$ decay, a
kink $1/|\xi|^2$, and a function with $k$ integrable derivatives decays
faster than $|\xi|^{-k}$. Sharp spike waveforms and stimulation artefacts
therefore spread power over a broad high-frequency range, which is why
they contaminate gamma and high-frequency band estimates.

\section{Autoregressive spectra and Yule--Walker}
For a stationary AR($p$) process
$X_t=\sum_{k=1}^pa_kX_{t-k}+\varepsilon_t$ with white innovations of
variance $\sigma^2$, multiplying by $X_{t-m}$ ($m\ge1$) and taking
expectations gives Result 67:
\[
\gamma(m)=\sum_{k=1}^pa_k\gamma(m-k),\qquad m=1,\ldots,p,
\]
and $\gamma(0)=\sum_ka_k\gamma(k)+\sigma^2$. Solving these Toeplitz
equations estimates the coefficients from sample autocovariances. The
spectral density is
\[
S(f)=\frac{\sigma^2/f_s}{\bigl|1-\sum_{k=1}^pa_ke^{-2\pi ikf/f_s}\bigr|^2}.
\]
\begin{example}[AR(2) resonance]
For $X_t=2r\cos\theta\,X_{t-1}-r^2X_{t-2}+\varepsilon_t$ with $0<r<1$, the
characteristic polynomial has roots $re^{\pm i\theta}$. The spectrum
has a peak near $f_0=\theta f_s/(2\pi)$ whose width is approximately
$(1-r)f_s/\pi$, and the autocorrelation is a damped cosine
$\propto r^{|k|}\cos(\theta|k|+\psi)$. With $f_s=250$\,Hz, $\theta$
chosen for $8$\,Hz and $r=0.98$, the process looks like a waxing and
waning theta rhythm with a $1.6$\,Hz-wide spectral peak, yet it is a
linear stochastic resonance: there is no limit cycle, and amplitude and
phase are driven entirely by noise. AR(2) processes are the standard
null model for ``is this oscillation more than filtered noise?''.
\end{example}

\section{Wiener--Hopf prediction and decoding}
The best linear predictor of $Y_t$ from $X_{t-k}$, $k=0,\ldots,K$, is
$\hat Y_t=\sum_kw_kX_{t-k}$ with weights solving the normal equations
(Result 68)
\[
\sum_{k=0}^K w_k\,\gamma_{xx}(j-k)=\gamma_{yx}(j),\qquad j=0,\ldots,K.
\]
\begin{proof}
The squared error is a convex quadratic in $w$. Setting its derivative
with respect to $w_j$ to zero is equivalent to
$\E[(Y_t-\hat Y_t)X_{t-j}]=0$: the residual is orthogonal to every
regressor, which is the orthogonal projection theorem (Result 26).
\end{proof}
With unlimited lags the solution in frequency is the Wiener filter
$W(f)=S_{yx}(f)/S_{xx}(f)$ (non-causal) and the explained fraction of
variance at frequency $f$ is the coherence $C_{yx}(f)$. Coherence is
therefore the fraction of $Y$'s power at $f$ that is linearly
predictable from $X$ \emph{at all lags, past and future}: it is a
symmetric measure of linear dependence, not of directed influence.

\section{Shift, modulation, and the meaning of phase}
By Result 60, $x(t-\tau)$ has transform $e^{-2\pi if\tau}X(f)$. A pure
delay shows up as a phase that is linear in frequency with slope
$-2\pi\tau$. Conversely, if the cross-spectrum phase between two
channels is approximately linear in $f$ over a band, its slope gives a
delay estimate (the \emph{phase-slope index} is a robust version of this
idea). A constant phase offset, independent of frequency, does not
correspond to a time delay at all; it is what a reference subtraction,
a derivative, or a Hilbert transform produces.

\section*{Exercises for this chapter}
\begin{exercise}
Prove Parseval's identity for Fourier series on $[0,T]$ from the
orthogonality relations of Result 57 for trigonometric polynomials.
\end{exercise}
\begin{exercise}
Derive the Yule--Walker equations for AR(1) and show that $\gamma(k)=a^{|k|}\sigma^2/(1-a^2)$.
\end{exercise}
\begin{exercise}
For the AR(2) example, find $\theta$ for an 8\,Hz peak at $f_s=250$\,Hz
and compute the coefficients.
\end{exercise}
\begin{exercise}
Show that $y(t)=x(t-\tau)$ has coherence one with $x$ at every frequency
and cross-spectral phase $-2\pi f\tau$.
\end{exercise}
